\documentclass[12pt]{article}
\usepackage[utf8]{inputenc}
\usepackage[T1]{fontenc}
\usepackage[margin=1in]{geometry}
\usepackage{setspace}
\usepackage{xcolor}
\usepackage[colorlinks=true, urlcolor=blue!60!black, linkcolor=blue!60!black]{hyperref}
\setlength{\parindent}{2em}
\setlength{\parskip}{1.25em}
\renewcommand{\baselinestretch}{1.0}
\begin{document}
\begin{center}
{\LARGE \textit{Katzenbach v. McClung}, 379 U.S. 294 (1964)\par}
\vspace{0.8em}
{\large Case Brief\par}
\end{center}
\vspace{1em}

\section{Memory Jogger}

In \href{https://www.courtlistener.com/opinion/106951/katzenbach-v-mcclung/}{\textit{Katzenbach v. McClung}}, the Supreme Court unanimously held that Congress may use the Commerce Clause to ban racial discrimination at a purely local restaurant---Ollie's Barbecue---so long as a substantial portion of its food has moved in interstate commerce and Congress had a rational basis to conclude that such discrimination, in the aggregate, burdens interstate trade.

\section{Detailed Case Facts}

Ollie McClung, Sr., and Ollie McClung, Jr., owned and operated Ollie's Barbecue, a family-owned restaurant in Birmingham, Alabama, that specialized in barbecued meats and homemade pies and had a seating capacity of 220 customers (\href{https://supreme.justia.com/cases/federal/us/379/294/}{379 U.S. at 296}). The restaurant sat on a state highway 11 blocks from an interstate highway and a somewhat greater distance from railroad and bus stations (\href{https://supreme.justia.com/cases/federal/us/379/294/}{379 U.S. at 296}). It catered to a local family and white-collar clientele, employed 36 people (two-thirds of whom were African American), and since opening in 1927 had refused to serve African Americans in its dining room while offering them only take-out service (\href{https://supreme.justia.com/cases/federal/us/379/294/}{379 U.S. at 296--297}).

In the 12 months preceding the enactment of the Civil Rights Act of 1964, Ollie's Barbecue purchased approximately \$150,000 worth of food locally; of that total, \$69,683---or 46\%---consisted of meat bought from a local supplier who had procured it from outside Alabama (\href{https://supreme.justia.com/cases/federal/us/379/294/}{379 U.S. at 296}). Title II of the Civil Rights Act of 1964, \S 201(a), 78 Stat. 243, guaranteed all persons the full and equal enjoyment of places of public accommodation without discrimination or segregation based on race, color, religion, or national origin (\href{https://supreme.justia.com/cases/federal/us/379/294/}{379 U.S. at 298}). Under \S\S 201(b)(2) and 201(c)(2), a restaurant principally engaged in selling food for consumption on the premises was covered if ``it serves or offers to serve interstate travelers or a substantial portion of the food which it serves . . . has moved in commerce'' (\href{https://supreme.justia.com/cases/federal/us/379/294/}{379 U.S. at 298}).

Ollie's Barbecue conceded that it fell within the statutory text and had been operating in violation of the Act since July 2, 1964 (\href{https://supreme.justia.com/cases/federal/us/379/294/}{379 U.S. at 297--298}). Unlike the companion case \href{https://supreme.justia.com/cases/federal/us/379/241/}{\textit{Heart of Atlanta Motel, Inc. v. United States}}, 379 U.S. 241 (1964), the United States made no claim that interstate travelers frequented Ollie's Barbecue or that the State of Alabama supported its segregation (\href{https://supreme.justia.com/cases/federal/us/379/294/}{379 U.S. at 298}). The sole constitutional issue was whether Congress could regulate a local restaurant solely because roughly \$70,000 of the food it served annually had moved in interstate commerce.

\section{Procedural History}

The McClungs sued Acting Attorney General Nicholas Katzenbach and a United States Attorney in the United States District Court for the Northern District of Alabama, seeking an injunction barring enforcement of Title II against Ollie's Barbecue (\href{https://supreme.justia.com/cases/federal/us/379/294/}{379 U.S. at 295}; \href{https://clearinghouse.net/case/11136/}{Civil Rights Litigation Clearinghouse}). The federal defendants moved to dismiss for want of equity jurisdiction because Title II authorized only preventive relief and no enforcement action had been threatened (\href{https://supreme.justia.com/cases/federal/us/379/294/}{379 U.S. at 295}).

On September 17, 1964, a three-judge District Court rejected the motion to dismiss, found that a substantial portion of the restaurant's food had moved in interstate commerce and that compliance would cause Ollie's Barbecue to lose a substantial amount of business, and held Title II unconstitutional as applied to the restaurant, issuing an injunction against its enforcement (\textit{McClung v. Katzenbach}, 233 F. Supp. 815 (N.D. Ala. 1964); \href{https://supreme.justia.com/cases/federal/us/379/294/}{379 U.S. at 295, 297}; \href{https://clearinghouse.net/case/11136/}{Civil Rights Litigation Clearinghouse}). The District Court reasoned that Congress had impermissibly legislated a conclusive presumption without a demonstrable connection between out-of-state food purchases and restaurant discrimination (\href{https://supreme.justia.com/cases/federal/us/379/294/}{379 U.S. at 297}).

The government took a direct appeal to the Supreme Court under 28 U.S.C. \S\S 1252 and 1253; Justice Black stayed the injunction pending appeal, and the Court noted probable jurisdiction at 379 U.S. 802 (\href{https://supreme.justia.com/cases/federal/us/379/294/}{379 U.S. at 295}; \href{https://clearinghouse.net/case/11136/}{Civil Rights Litigation Clearinghouse}). Argued October 5, 1964, alongside \href{https://supreme.justia.com/cases/federal/us/379/241/}{\textit{Heart of Atlanta Motel}}, the case was decided on December 14, 1964, reversing the District Court (\href{https://supreme.justia.com/cases/federal/us/379/294/}{379 U.S. at 294--295, 305}).

\section{Judicial Votes}

The Supreme Court reversed the District Court by a unanimous vote of \textbf{9--0} (\href{https://supreme.justia.com/cases/federal/us/379/294/}{379 U.S. at 294--305}):

\begin{itemize}
\item \textbf{Majority Opinion:} Justice Tom C. Clark authored the opinion of the Court, joined in full by Chief Justice Earl Warren and Justices John Marshall Harlan II, William J. Brennan, Jr., Potter Stewart, Byron R. White, William O. Douglas, and Arthur J. Goldberg (\href{https://supreme.justia.com/cases/federal/us/379/294/}{379 U.S. at 295}; \href{https://supreme.justia.com/cases/federal/us/379/241/}{379 U.S. at 279, 291}).
\item \textbf{Separate Concurrences:} Justices Hugo L. Black, William O. Douglas, and Arthur J. Goldberg each filed a separate concurring opinion applicable to both \href{https://supreme.justia.com/cases/federal/us/379/241/}{\textit{Heart of Atlanta Motel}} and \href{https://www.courtlistener.com/opinion/106951/katzenbach-v-mcclung/}{\textit{Katzenbach v. McClung}}, printed in \textit{Heart of Atlanta Motel} at 379 U.S. 268 (Black, J.), 379 U.S. 279 (Douglas, J.), and 379 U.S. 291 (Goldberg, J.) (\href{https://supreme.justia.com/cases/federal/us/379/241/}{379 U.S. at 268--293}).
\item \textbf{Dissents:} None.
\end{itemize}

\section{Holding}

\textbf{Judgment:} Reversed. Title II of the Civil Rights Act of 1964, \S\S 201(a), (b)(2), and (c)(2), is constitutional under the Commerce Clause (U.S. Const. art. I, \S 8, cl. 3) and the Necessary and Proper Clause (art. I, \S 8, cl. 18) as applied to a local restaurant serving food a substantial portion of which has moved in interstate commerce (\href{https://supreme.justia.com/cases/federal/us/379/294/}{379 U.S. at 301--305}).

\textbf{Rule Announced:} Congress may regulate local commercial establishments under the Commerce Clause when, in light of the facts and testimony before it, Congress has a \textit{rational basis} for finding that the regulated class of activities, considered in the aggregate, exerts a substantial effect on interstate commerce and that the chosen regulatory scheme is a reasonable and appropriate means to protect commerce (\href{https://supreme.justia.com/cases/federal/us/379/294/}{379 U.S. at 300--304}). Neither formal statutory findings nor case-by-case proof that an individual business affects interstate commerce is constitutionally required (\href{https://supreme.justia.com/cases/federal/us/379/294/}{379 U.S. at 302--304}).

\textbf{Holding vs. Dicta:} The Court's core holding is that Title II validly applies to a restaurant purchasing a substantial portion of its food from out of state, even without proof that it serves interstate travelers, and that the absence of formal legislative findings or individualized nexus determinations does not invalidate the statute (\href{https://supreme.justia.com/cases/federal/us/379/294/}{379 U.S. at 303--305}). The Court also made a narrow procedural holding treating the complaint as a declaratory judgment action under 28 U.S.C. \S\S 2201--2202 because of the urgent need to resolve Title II's validity alongside \textit{Heart of Atlanta Motel}, while admonishing in dicta that courts should ordinarily withhold declaratory relief where a statute provides a special enforcement proceeding (\href{https://supreme.justia.com/cases/federal/us/379/294/}{379 U.S. at 295--296}). Similarly, Justice Clark's sweeping language quoting \textit{Gibbons v. Ogden}, 9 Wheat. 1, 195 (1824), and describing congressional commerce power as ``broad and sweeping'' functioned as expansive framing later cabined by the Rehnquist Court in \href{https://supreme.justia.com/cases/federal/us/514/549/}{\textit{United States v. Lopez}}, 514 U.S. 549 (1995), and \textit{United States v. Morrison}, 529 U.S. 598 (2000), which limited aggregation to \textit{economic} activities (\href{https://www.everycrsreport.com/reports/RL32844.html}{CRS Report RL32844}).

\section{Analysis of Opinions}

\subsection{Majority Opinion (Justice Clark)}

Justice Clark first disposed of the government's procedural objection by construing the complaint as a declaratory judgment action so the Court could promptly settle Title II's constitutionality alongside \href{https://supreme.justia.com/cases/federal/us/379/241/}{\textit{Heart of Atlanta Motel}} (\href{https://supreme.justia.com/cases/federal/us/379/294/}{379 U.S. at 295--296}). Turning to the merits, he noted that \textit{Heart of Atlanta Motel} already foreclosed the appellees' Fifth, Ninth, Tenth, and Thirteenth Amendment claims as well as their reliance on the \textit{Civil Rights Cases}, 109 U.S. 3 (1883) (\href{https://supreme.justia.com/cases/federal/us/379/294/}{379 U.S. at 298 \& n. 1}).

Reviewing the prolonged congressional hearings, Justice Clark highlighted four evidentiary links between restaurant segregation and interstate commerce: (1) per capita spending by African Americans in restaurants was depressed in segregated areas, causing restaurants to sell less food and consequently purchase less food from out of state; (2) discrimination imposed an artificial restriction on the market and depressed general business conditions through community unrest; (3) denying dining-room service directly obstructed interstate travel by African Americans, ``for one can hardly travel without eating''; and (4) segregation deterred skilled professionals and new industries from relocating to segregated regions (\href{https://supreme.justia.com/cases/federal/us/379/294/}{379 U.S. at 299--300}).

Although Ollie's Barbecue's individual food purchases were trivial compared to national commerce, Justice Clark invoked the aggregation principle of \textit{Wickard v. Filburn}, 317 U.S. 111, 127--128 (1942): an individual actor's trivial contribution does not remove it from federal regulation when, taken together with many others similarly situated, the cumulative impact is far from trivial (\href{https://supreme.justia.com/cases/federal/us/379/294/}{379 U.S. at 300--301}). Rejecting the argument that interstate commerce ends once goods come to rest in-state, he explained that such ``come to rest'' cases governed state taxation, not federal regulatory power (\href{https://supreme.justia.com/cases/federal/us/379/294/}{379 U.S. at 302}). Relying on \textit{United States v. Darby}, 312 U.S. 100 (1941), and \textit{United States v. Carolene Products Co.}, 304 U.S. 144 (1938), he concluded that Congress need neither include formal statutory findings nor provide case-by-case proof of commercial impact: once the Court finds a rational basis for Congress's determination that the regulated class affects commerce, ``our investigation is at an end'' (\href{https://supreme.justia.com/cases/federal/us/379/294/}{379 U.S. at 303--304}).

\subsection{Concurring Opinion (Justice Black)}

Justice Black concurred in upholding Title II under the Commerce Clause and the Necessary and Proper Clause, emphasizing that 46\% (\$69,683) of Ollie's Barbecue's food supplies was plainly a ``substantial'' portion of interstate goods and that aggregate racial exclusion of roughly 20 million African American consumers could distort interstate food markets and deter highway travel (\href{https://supreme.justia.com/cases/federal/us/379/241/}{379 U.S. at 274--276}). Crucially, however, Justice Black sounded a federalism warning that foreshadowed the Rehnquist Court: whether local operations sufficiently affect interstate commerce is ``ultimately a judicial, rather than a legislative, question,'' and ``every remote, possible, speculative effect on commerce should not be accepted as an adequate constitutional ground to uproot and throw into the discard all our traditional distinctions between what is purely local . . . and what affects the national interest'' (\href{https://supreme.justia.com/cases/federal/us/379/241/}{379 U.S. at 273, 275}). He noted that an isolated, remote lunchroom selling only to locals and buying almost all supplies locally might remain beyond Congress's commerce power (\href{https://supreme.justia.com/cases/federal/us/379/241/}{379 U.S. at 275}). Because the Commerce Clause was sufficient, he saw no need to reach \S 5 of the Fourteenth Amendment (\href{https://supreme.justia.com/cases/federal/us/379/241/}{379 U.S. at 279}).

\subsection{Concurring Opinion (Justice Douglas)}

While joining the Court's opinions, Justice Douglas expressed reluctance to rest civil rights legislation solely on the Commerce Clause, arguing that the right to be free of racial discrimination ``occupies a more protected position in our constitutional system than does the movement of cattle, fruit, steel and coal across state lines'' (\href{https://supreme.justia.com/cases/federal/us/379/241/}{379 U.S. at 279}). He preferred to ground Title II in \S 5 of the Fourteenth Amendment, reasoning under \textit{Shelley v. Kraemer}, 334 U.S. 1 (1948), that state judicial enforcement of trespass laws against African American patrons constituted state action (\href{https://supreme.justia.com/cases/federal/us/379/241/}{379 U.S. at 280--283}). Resting on the Fourteenth Amendment, he argued, would have a ``more settling effect,'' eliminating endless litigation over whether a specific restaurant purchased enough interstate food or served interstate travelers (\href{https://supreme.justia.com/cases/federal/us/379/241/}{379 U.S. at 280, 286}).

\subsection{Concurring Opinion (Justice Goldberg)}

Justice Goldberg joined the Court's opinions and judgments under the Commerce Clause, but wrote separately to underscore that ``[t]he primary purpose of the Civil Rights Act of 1964 . . . is the vindication of human dignity, and not mere economics'' (\href{https://supreme.justia.com/cases/federal/us/379/241/}{379 U.S. at 291}). Quoting Senate Report No. 872, he emphasized that discrimination inflicts personal humiliation rather than merely affecting ``dollars and cents, hamburgers and movies'' (\href{https://supreme.justia.com/cases/federal/us/379/241/}{379 U.S. at 292}). Reaffirming his concurrence in \textit{Bell v. Maryland}, 378 U.S. 226, 317 (1964), he concluded that Congress possessed ample authority under \textit{both} \S 5 of the Fourteenth Amendment and the Commerce Clause (\href{https://supreme.justia.com/cases/federal/us/379/241/}{379 U.S. at 292--293}).

\section{Examples: Future Applications}

\begin{enumerate}
\item \textbf{Same-Side (Valid under \textit{McClung} and Post-\textit{Lopez} Doctrine):} Congress prohibits racial discrimination by local neighborhood grocery stores and coffee shops that purchase at least 25\% of their inventory or coffee beans from out-of-state distributors. \textit{Why it falls on the same side:} Like Ollie's Barbecue, these retail establishments engage in quintessentially commercial, economic activity and satisfy an express jurisdictional element tying their operations to goods that moved in interstate commerce; aggregating discriminatory exclusions across such stores rationally demonstrates a substantial effect on interstate food markets (\href{https://supreme.justia.com/cases/federal/us/379/294/}{379 U.S. at 300--304}; \href{https://www.everycrsreport.com/reports/RL32844.html}{CRS Report RL32844}).

\item \textbf{Same-Side (Valid under \textit{McClung} and \textit{Gonzales v. Raich}):} Congress enforces a comprehensive federal statute regulating the national market for prescription medications against a small, intrastate compounding pharmacy that sells only to local residents using active ingredients shipped from another state. \textit{Why it falls on the same side:} Compounding and selling pharmaceuticals is economic activity involving interstate inputs, and exempting local retail actors would rationally undercut the broader interstate regulatory scheme, matching both \textit{McClung}'s aggregation logic and \href{https://supreme.justia.com/cases/federal/us/545/1/}{\textit{Gonzales v. Raich}}, 545 U.S. 1 (2005).

\item \textbf{Opposite-Side (Invalid under Rehnquist Court Limits):} Congress enacts a statute making it a federal crime to possess a knife within 1,000 feet of a public park, with no jurisdictional element requiring the knife to have crossed state lines, arguing that violent crime in parks deters interstate tourism and reduces national productivity. \textit{Why it falls on the opposite side:} Under \href{https://supreme.justia.com/cases/federal/us/514/549/}{\textit{United States v. Lopez}}, 514 U.S. 549 (1995), mere possession of a weapon is non-economic criminal conduct; unlike Ollie's Barbecue's commercial food sales, allowing Congress to aggregate non-economic conduct through an attenuated ``costs of crime'' chain would convert the Commerce Clause into a general federal police power (\href{https://www.everycrsreport.com/reports/RL32844.html}{CRS Report RL32844}).

\item \textbf{Opposite-Side (Invalid under \textit{Morrison} and \textit{NFIB v. Sebelius}):} Congress passes a federal law compelling individuals who do not currently dine at restaurants to purchase at least two restaurant meals per month to stimulate interstate agricultural demand. \textit{Why it falls on the opposite side:} Even though restaurant dining is commercial in the aggregate under \textit{McClung}, \href{https://supreme.justia.com/cases/federal/us/567/519/}{\textit{National Federation of Independent Business v. Sebelius}}, 567 U.S. 519 (2012), holds that the Commerce Clause empowers Congress to regulate existing commercial \textit{activity}, not to compel inactive individuals into commerce.

\item \textbf{Fence-Sitter (Borderline Application):} Congress invokes the Commerce Clause to ban discrimination by a tiny, rural farm-to-table diner that sources 99\% of its ingredients from its own backyard garden and neighboring in-state farms, serves exclusively local residents, and uses only an out-of-state brand of salt and a commercial stove manufactured in another state. \textit{Why it is a fence-sitter:} On one hand, the diner conducts commercial retail activity and uses minor out-of-state inputs, and \textit{McClung}'s deferential rational-basis review gives Congress broad leeway over a class of economic establishments. On the other hand, Justice Black's concurrence explicitly cautioned that an ``isolated and remote lunchroom'' buying almost all supplies locally might lie beyond federal reach (\href{https://supreme.justia.com/cases/federal/us/379/241/}{379 U.S. at 275}), and Title II itself requires that a ``substantial portion'' of the food served have moved in interstate commerce (\href{https://supreme.justia.com/cases/federal/us/379/294/}{379 U.S. at 298}).
\end{enumerate}

\section{Critique}

\textbf{Originalist and Textualist Perspectives:} Originalist scholars and jurists have long viewed \href{https://www.courtlistener.com/opinion/106951/katzenbach-v-mcclung/}{\textit{Katzenbach v. McClung}} as the high-water mark of an unmoored Commerce Clause that erased the Founding-era distinction between interstate trade and local police power. Scholars such as Richard A. Epstein and Raoul Berger have argued that the original public meaning of ``Commerce . . . among the several States'' was limited to the interstate trade, sale, and transportation of goods across state lines, not purely intrastate retail service occurring after goods have come to rest (\href{https://lawcat.berkeley.edu/record/1118252/files/fulltext.pdf}{Berkeley Law Commentary on \textit{Lopez} and \textit{Morrison}}; \href{https://constitution.heritage.org/essays/a1-s8-c3-b/}{Heritage Guide to the Constitution}). Similarly, Justice Clarence Thomas has repeatedly criticized the ``substantial effects'' and aggregation doctrines---epitomized by \textit{Wickard} and \textit{McClung}---for permitting Congress to appropriate state police powers over local conduct (\href{https://dspace.loyno.edu/jspui/bitstream/123456789/131/1/Article\%204\%20-\%20McGoldrick_PRINT.pdf}{Loyola Law Review}; \href{https://courses.lumenlearning.com/clinton-buslegalenv/chapter/4-2-the-commerce-clause/}{Lumen Learning}). Critics of \textit{McClung}'s rational-basis deference also warn that treating any local establishment using out-of-state ingredients as affecting commerce leaves virtually no judicially enforceable boundary on federal power (\href{https://digitalcommons.law.byu.edu/cgi/viewcontent.cgi?article=1252&context=jpl}{BYU Journal of Public Law}). In the Rehnquist Court era, \href{https://supreme.justia.com/cases/federal/us/514/549/}{\textit{Lopez}} (1995) and \textit{Morrison} (2000) preserved \textit{McClung} as precedent---because running a restaurant is economic activity and Title II contained an express jurisdictional hook---while rejecting its unbounded rational-basis rhetoric for non-economic conduct (\href{https://www.everycrsreport.com/reports/RL32844.html}{CRS Report RL32844}).

\textbf{Progressive and Doctrinal Perspectives:} Progressive and civil-rights scholars applaud the outcome of \textit{McClung} for dismantling Jim Crow in public accommodations, yet many echo Justices Douglas and Goldberg in criticizing the Court's doctrinal choice to rely solely on the Commerce Clause rather than \S 5 of the Fourteenth Amendment (\href{https://supreme.justia.com/cases/federal/us/379/241/}{379 U.S. at 279--280, 291--293}; \href{https://constitution.congress.gov/browse/essay/artI-S8-C3-6-8/ALDE_00013425/}{Constitution Annotated}). Treating racial exclusion as a problem of interstate meat shipments rather than equal citizenship and human dignity risked trivializing a moral imperative and left future civil rights statutes---such as the Violence Against Women Act invalidated in \textit{Morrison}---vulnerable when the Rehnquist Court narrowed both the Commerce Clause and the \textit{Civil Rights Cases}' state-action barrier (\href{https://biotech.law.lsu.edu/cases/preemption/us_v_morrison.htm}{\textit{United States v. Morrison}}; \href{https://scholarship.law.vanderbilt.edu/cgi/viewcontent.cgi?article=1531&context=vlr}{Vanderbilt Law Review}). At the same time, pragmatic doctrinal defenders note that grounding Title II in the Commerce Clause secured a unanimous 9--0 Court, avoided a divisive battle over overruling the \textit{Civil Rights Cases}, and ensured rapid nationwide compliance (\href{https://supreme.justia.com/cases/federal/us/379/294/}{Justia Commentary}).

\textbf{Underlying Competing Values:} \textit{Katzenbach v. McClung} crystallizes the enduring constitutional tension between \textit{enumerated-powers federalism} (preserving a sphere of local self-governance and private commercial autonomy) and \textit{national integration and equal dignity} (empowering the federal legislature to dismantle entrenched regional caste systems in an interconnected national economy).

\section{Key Quotations}

\begin{quote}
``Of course, the mere fact that Congress has said when particular activity shall be deemed to affect commerce does not preclude further examination by this Court. But where we find that the legislators, in light of the facts and testimony before them, have a rational basis for finding a chosen regulatory scheme necessary to the protection of commerce, our investigation is at an end.''
\end{quote}
\noindent --- Justice Tom C. Clark, Opinion of the Court, 379 U.S. at 303--304 (\href{https://supreme.justia.com/cases/federal/us/379/294/}{Justia} \textbar{} \href{https://www.courtlistener.com/opinion/106951/katzenbach-v-mcclung/}{CourtListener}).

\begin{quote}
``We believe that this testimony afforded ample basis for the conclusion that established restaurants in such areas sold less interstate goods because of the discrimination, that interstate travel was obstructed directly by it, that business in general suffered, and that many new businesses refrained from establishing there as a result of it.''
\end{quote}
\noindent --- Justice Tom C. Clark, Opinion of the Court, 379 U.S. at 300 (\href{https://supreme.justia.com/cases/federal/us/379/294/}{Justia} \textbar{} \href{https://www.courtlistener.com/opinion/106951/katzenbach-v-mcclung/}{CourtListener}).

\begin{quote}
``The power of Congress in this field is broad and sweeping; where it keeps within its sphere and violates no express constitutional limitation it has been the rule of this Court, going back almost to the founding days of the Republic, not to interfere.''
\end{quote}
\noindent --- Justice Tom C. Clark, Opinion of the Court, 379 U.S. at 305 (\href{https://supreme.justia.com/cases/federal/us/379/294/}{Justia} \textbar{} \href{https://www.courtlistener.com/opinion/106951/katzenbach-v-mcclung/}{CourtListener}).

\begin{quote}
``A decision based on the Fourteenth Amendment would have a more settling effect, making unnecessary litigation over whether a particular restaurant or inn is within the commerce definitions of the Act or whether a particular customer is an interstate traveler.''
\end{quote}
\noindent --- Justice William O. Douglas, Concurring Opinion (applicable also to \textit{McClung}), \textit{Heart of Atlanta Motel, Inc. v. United States}, 379 U.S.\ 241, 280 (1964) (\href{https://supreme.justia.com/cases/federal/us/379/241/}{Justia}).

\begin{quote}
``The primary purpose of the Civil Rights Act of 1964, however, as the Court recognizes, and as I would underscore, is the vindication of human dignity, and not mere economics.''
\end{quote}
\noindent --- Justice Arthur J. Goldberg, Concurring Opinion (applicable also to \textit{McClung}), \textit{Heart of Atlanta Motel, Inc. v. United States}, 379 U.S.\ 241, 291 (1964) (\href{https://supreme.justia.com/cases/federal/us/379/241/}{Justia}).

\end{document}
